\section{符号说明}

正文所用主要符号如表~\ref{tab:symbols} 所示。

\begin{table}[htbp]
\centering
\caption{主要符号及其含义}
\label{tab:symbols}
\begin{tabular}{cll}
\toprule
符号 & 含义 & 单位或取值范围 \\
\midrule
$\mathcal I,\mathcal R,\mathcal T$ & 任务、区域与小时集合 & 集合 \\
$\mathcal I^{\mathrm{flex}}$ & 批量推理与 AI 训练任务集合 & $\mathcal I$ 的子集 \\
$i,r,t,s,k$ & 任务、执行区域、小时、开工时刻与任务类型索引 & 无量纲 \\
$o_i,k_i$ & 任务 $i$ 的来源区域与任务类型 & 类别 \\
$a_i,d_i,f_i^{\mathrm{late}}$ & 到达时刻、连续执行时长与最晚完成时点 & h \\
$g_i$ & 任务 $i$ 持续占用的等效 GPU 数 & 等效 GPU \\
$\ell_i^{\max},\ell_{or}$ & 最大允许时延与来源区至执行区单向时延 & ms \\
$x_{irs}$ & 任务 $i$ 是否在区域 $r$、时刻 $s$ 开工 & $\{0,1\}$ \\
$\omega_{it}(s)$ & 任务 $i$ 在小时 $t$ 内的实际重叠时长 & h \\
$\widehat D_{rkt}$ & 区域—类型逐小时 GPU 需求预测值 & 等效 GPU \\
$G_r^{\max}$ & 区域 $r$ 的可调度 GPU 容量 & 等效 GPU \\
$p_k^{GPU}$ & 类型 $k$ 的单位等效 GPU 平均 IT 功率 & MW/GPU \\
$P_{rt}^{AI},P_{rt}^{NA},P_{rt}^{BAI}$ & 调度 AI、固定非 AI 与基线 AI IT 功率 & MW \\
$P_{rt}^{IT},P_{rt}^{fac}$ & 总 IT 功率与设施侧总功率 & MW \\
$P_r^{IT,\max},P_r^{fac,\max}$ & IT 功率与设施功率上限 & MW \\
$\mathrm{PUE}_r$ & 区域 $r$ 的能源使用效率 & 无量纲 \\
$R_{rt}^{av},R_{rt}^{use}$ & 可用新能源与直接供负荷的新能源功率 & MW \\
$C_{rt}^{RE},C_{rt}^{G},C_{rt}$ & 新能源充电、电网充电与总充电功率 & MW \\
$D_{rt}$ & 储能放电功率 & MW \\
$B_{rt}$ & 电网购电功率 & MW \\
$S_{rt}$ & 新能源外送或售电功率 & MW \\
$K_{rt}$ & 弃风弃光功率 & MW \\
$E_{rt},E_r^0$ & 时段末绝对 SOC 与初始 SOC & MWh \\
$E_r^{\min},E_r^{\max}$ & 储能 SOC 下限与容量上限 & MWh \\
$C_r^{\max},D_r^{\max}$ & 最大充电与放电功率 & MW \\
$\eta_r^c,\eta_r^d$ & 储能充电效率与放电效率 & 无量纲 \\
$P_{rt}^{grid,\max}$ & 区域购电功率上限 & MW \\
$P_{rt}^{export,\max},S_r^{\max}$ & 区域外送边界与储能表售电上限 & MW \\
$N_{rt}$ & 净购电功率，即 $B_{rt}-S_{rt}$ & MW \\
$H_r$ & 区域峰值正向净购电功率 & MW \\
$V_r$ & 区域净购电波动指标 & MW \\
$C_{\mathrm{op}}$ & 购电支出扣除售电收入后的运行成本 & 元 \\
$E_{\mathrm{CO2}}$ & 电网购电产生的累计碳排放 & tCO2 \\
$L_{\mathrm{net}}$ & GPU-hour 加权网络时延指标 & ms \\
$U_{\mathrm{RE}}$ & 新能源利用率 & \% \\
$Q_{\mathrm{service}}$ & 基于题内可观测量的服务质量代理指标 & \%或无量纲 \\
\bottomrule
\end{tabular}
\end{table}

表中 $B$、$R$、$C$、$D$、$S$、$K$ 和 $E$ 分别对应购电、新能源、充电、放电、外送、弃电及储能状态，是能源平衡与 SOC 递推的核心量。符号 $P$ 统一表示功率，利用上下标区分 AI、非 AI、设施负荷及不同接入边界，从而避免任务负荷与电力流混淆。目标指标 $C_{\mathrm{op}}$、$E_{\mathrm{CO2}}$、$L_{\mathrm{net}}$、$U_{\mathrm{RE}}$、$Q_{\mathrm{service}}$、$H_r$ 和 $V_r$ 均由同一完整调度候选重新计算，以保证不同问题和情景之间的评价口径一致。
